\documentclass[../main.tex]{subfiles}

\begin{document}
\vspace{-0.5ex}
    \subsection{Rental PS Rijal}
\begin{lstlisting}
 #include <iostream>
 using namespace std;
\end{lstlisting}
    \ilz{}{\Script} diatas digunakan untuk memberikan kemampuan untuk menggunakan fungsi-fungsi dan objek-objek yang tersedia dalah {\header} {\file} \inputscript{iostream}. {\Script} pada baris kedua \inputscript{using namespace std;} memberikan kemampuan untuk menggunakan fungsi-fungsi dan objek-objek dari \textit{standard library} tanpa menulis \inputscript{std::}.

\begin{lstlisting}
 const int user_max = 20;
 string users[user_max * 2] = { [0] = "admin", [1] = users[0] };
 bool user_membership[user_max];
 int user_bills[user_max];
 int user_logged_id = -1;
 int user_count = 1;
 int console_prices[4] = { 0, 5000, 10000, 15000 };
 const int food_max = 20;
 string foods[food_max];
 int food_prices[food_max];
 int food_count = 0;
 int app_state = 0;
\end{lstlisting}
    \ilz{}{\Script} diatas digunakan untuk menginisialisasikan semua variabel yang dibutuhkan untuk menjalankan program Rental PS Rijal. Semua variabel yang diawali oleh kata \textit{user} merupakan variabel yang digunakan untuk menyimpan data pengguna atau berkaitan dengan data pengguna. Semua variabel yang diawali oleh kata \textit{food} merupakan variabel yang digunakan untuk menyimpan data makanan dan minuman. Variabel \inputscript{app\_state} merupakan variabel yang menyimpan data keadaan aplikasi.

\vspace{-0.15ex}
\begin{lstlisting}
 void pause() {
   cout << "Press any key to continue. . .\n";
   cin.clear(); cin.ignore(); cin.get(); }
\end{lstlisting}
    \ilz{}{\Script} diatas merupakan sebuah fungsi tanpa kembalian yang ditandakan oleh kata kunci \inputscript{void}. Fungsi tersebut akan menghentikan pengguna sejenak sampai dideteksi sebuah masukan apapun dari \textit{keyboard}.

\vspace{-0.15ex}
\begin{lstlisting}
 void input_invalid() {
   cin.clear(); cin.ignore();
   cout << "Masukan tidak valid\n"; }
\end{lstlisting}
    \ilz{}{\Script} diatas merupakan sebuah fungsi tanpa kembalian yang ditandakan oleh kata kunci \inputscript{void}. Fungsi tersebut akan memberitahu pengguna bahwa sebuah masukan tidak valid dan membersihkan masukan tidak valid tersebut dari {\ioi} {\stream}.

\vspace{-0.15ex}
\begin{lstlisting}
 int nselector(int l, int u, string c = ">> ") {
   int n;
   while (true) {
\end{lstlisting}
\begin{lstlisting}
     cout << c;
     cin >> n;
     if (cin.fail() || n < l || n > u) {
       input_invalid();
       continue; } break;
   } return n; }
\end{lstlisting}
    \ilz{}{\Script} diatas merupakan sebuah fungsi dengan kembalian berupa \inputscript{int} atau bilangan bulat. Fungsi tersebut akan menyaring masukan dari pengguna sesuai kedua parameter fungsi \inputscript{int l} dan \inputscript{int u}. Masukan yang dimasukan harus berada dalam lingkup kedua parameter tersebut untuk dianggap sebagai masukan yang valid.

\begin{lstlisting}
 bool login(string name, string password) {
   for (int i = 0; i < user_count; ++i) {
     if (name == users[i * 2]
       && password == users[i * 2 + 1]) {
       user_logged_id = i;
       return true; }
   } return false; }
\end{lstlisting}
    \ilz{}{\Script} diatas merupakan sebuah fungsi dengan kembalikan berupa tipe data \inputscript{bool} yang hanya mampu bernilai bernar atau salah. Fungsi tersebut menerima dua parameter yaitu nama pengguna dan kata sandi. Bila nama pengguna dan kata sandi cocok dengan data yang tersimpan maka keluaran fungsi akan benar dan salah untuk selainnya.

\begin{lstlisting}
 bool signin(string name, string password) {
   for (int i = 0; i < user_count; ++i) {
     if (name == users[i * 2]) return false;
   }
   users[user_count * 2] = name;
   users[user_count * 2 + 1] = password;
   user_membership[user_count] = false;
   user_bills[user_count] = 0;
   user_count++;
   return true; }
\end{lstlisting}
    \ilz{}{\Script} diatas merupakan sebuah fungsi dengan kembalikan berupa tipe data \inputscript{bool} yang hanya mampu bernilai bernar atau salah. Fungsi tersebut akan membuat akun pengguna baru bila nama pengguna belum pernah digunakan.

\begin{lstlisting}
 void topbar() {
   cout << "================================\n";
   cout << "=========Rental PS Rijal========\n";
   cout << "================================\n\n"; }
\end{lstlisting}
    \ilz{}{\Script} diatas merupakan sebuah fungsi tanpa kembalian yang ditandakan oleh kata kunci \inputscript{void}. Fungsi tersebut digunakan untuk mencetak {\header} atau \textit{bar} pada bagian atas program ketika program dijalankan.

\begin{lstlisting}
 void unlogged_page() {
   switch ((app_state & 0x6) >> 0x1) {
\end{lstlisting}
\begin{lstlisting}
     case 0: {
       topbar();
       cout << "1. Login\n2. Sign-In\n";
       cout << "3. Exit\n";
       app_state = (nselector(1, 3) << 0x1);
     } break;
     case 1: {
       cout << "================================\n";
       cout << "=============Log-In=============\n";
       cout << "================================\n\n";
       string name, password;
       cout << "Masukan Username: "; cin >> name;
       cout << "Masukan Password: "; cin >> password;
       if (app_state = login(name, password))
         cout << "Log-In Berhasil\n";
       else cout << "Log-In Tidak Berhasil\n";
       pause();
     } break;
     case 2: {
       cout << "================================\n";
       cout << "=============Sign-In============\n";
       cout << "================================\n\n";
       string name, password;
       cout << "Masukan Username: "; cin >> name;
       cout << "Masukan Password: "; cin >> password;
       if (signin(name, password)) cout << "Sing-In Berhasil\n";
       else cout << "Username sudah ada\n";
       pause(); app_state = 0;
     } break; } }
\end{lstlisting}
    \ilz{}{\Script} diatas mengandung sebuah fungsi yang digunakan untuk menggambarkan halaman ketika pengguna belum \textit{log-in} atau masuk. Halaman tersebut akan memberikan tiga pilihan yakni \textit{log-in}, \textit{sign-in} dan \textit{exit} yang akan bekerja sesuai namanya.

\begin{lstlisting}
 void logged_page() {
   topbar();
   switch ((app_state & 0x6) >> 0x1) {
     case 0: {
       cout << "1. Melakukan Rental PS\n";
       cout << "2. Memesan Makanan\n";
       cout << "3. Daftar Member\n";
       cout << "4. log-out\n";
       app_state |= (nselector(1, 4) << 0x1);
     } break;
     case 1: {
       cout << "Harga Rental PS\n";
       cout << "1. PS 3: 1 Jam (Rp. 5.000)\n";
       cout << "2. PS 4: 1 Jam (Rp. 10.000)\n";
       cout << "3. PS 5: 1 Jam (Rp. 15.000)\n";
       cout << "4. Kembali\n";
       {
         int n = nselector(1, 4);
         if (n == 4) break;
         int m = nselector(1, 9*9*9*9, "Jumlah Jam: ");
         float discount = (float)(console_prices[n] * m) *
           ((0.1f * (m >= 3 && m < 6)) + (0.2f * (m >= 6))
\end{lstlisting}
\begin{lstlisting}
            + (user_membership[user_logged_id] * 0.2f));
         user_bills[user_logged_id]
           += (console_prices[n] * m) - discount;
         pause(); }
       app_state = 1;
     } break;
     case 2: {
       cout << "List Makanan/Minuman\n";
       if (food_count == 0) {
         cout << "Makanan Belum Tersedia\n";
         app_state = 1; pause();
         break; }
       for (int i = 0, j = 0; i < food_count; ++i) {
         cout << i + 1 << ". "
            << foods[i] << " (Rp. "
            << food_prices[i] << ")\n";
       } cout << food_count + 1 << ". Kembali\n";
       int n = nselector(1, food_count + 1);
       if (n == food_count + 1) {
         app_state = 1;
         break; }
       cout << "Makanan yang dipilih: " << foods[n - 1] << "\n";
       cout << "Harga: " << food_prices[n - 1] << "\n";
       user_bills[user_logged_id] += food_prices[n - 1]
             * nselector(1, 9*9*9, "Jumlah: ");
       cout << foods[n - 1] << " berhasil dipesan\n";
       pause();
     } break;
     case 3: {
       char choice = 'y';
       while (true) {
         cout << "Apakah anda ingin mendaftar sebagai"
            << " member rental PS?(Y/T): ";
         cin >> choice;
         if (cin.fail() ||
           ((choice != 'Y' && choice != 'y') &&
            (choice != 'T' && choice != 't'))) {
           input_invalid();
           continue;
         } break; }
       if (choice == 'Y' || choice == 'y') {
         if (user_membership[user_logged_id])
           cout << "Anda sudah terdaftar sebagai member\n";
         else {
           cout << "Anda berhasil mendaftar sebagai member\n";
           user_membership[user_logged_id] = true; }
       } app_state = 1; pause();
     } break; } }
\end{lstlisting}
    \ilz{}{\Script} diatas mengandung sebuah fungsi yang tidak mengembalikan nilai apapun atau fungsi \inputscript{void}.Fungsi diatas digunakan untuk mengurus hal apapun yang berkaitan dengan fungsi yang dibutuhkan untuk pengguna yang sudah \textit{log-in}.

\begin{lstlisting}
 void admin_page() {
   topbar();
   switch ((app_state & 0x6) >> 1) {
\end{lstlisting}
\begin{lstlisting}
     case 0: {
       cout << "1. Jual Makanan\n";
       cout << "2. Data Penyewa\n";
       cout << "3. log-out\n";
       app_state |= (nselector(1, 3) << 0x1);
     } break;
     case 1: {
       cout << "List Makanan/Minuman\n";
       if (food_count == 0) {
         cout << "Makanan/Minuman Belum Tersedia\n";
       } else {
         for (int i = 0; i < food_count; ++i) {
           if (app_state >> 4 == 0) cout << "~ ";
           else cout << i + 1 << ". ";
           cout << foods[i] << " (Rp. "
              << food_prices[i] << ")\n";
         }
         if (app_state >> 4 != 0)
           cout << food_count + 1 << ". Kembali\n";
         else cout << "\n";
       }
       if ((app_state >> 4) == 0) {
         cout << "1. Masukan Makanan\n";
         cout << "2. Hapus Makanan\n";
         cout << "3. Kembali\n";
         int n = nselector(1, 3);
         if (n == 1) {
           cout << "\33[H\33[2J"; topbar();
           cout << "Masukan Makanan/Minuman: ";
           cin.clear(); cin.ignore();
           getline(cin, foods[food_count]);
           cout << "Masukan Harga: ";
           cin >> food_prices[food_count];
           cout << "Makanan/Minuman Berhasil"
              << " Dimasukan\n";
           food_count++;
           pause(); app_state = 3; break;
         } else if (n == 3) {
           app_state = 1; break;
         }
         app_state |= (0x1 << 4);
         break;
       }
       int n = nselector(1, food_count + 1);
       if (n == food_count + 1) {
         app_state = 3;
         break;
       }
       cout << foods[n - 1] << " berhasil dihapus\n";
       for (int i = (n - 1); i < food_max - 1; ++i) {
         foods[i] = foods[i + 1];
         food_prices[i] = food_prices[i + 1];
       } food_count--;
       pause(); app_state = 3;
     } break;
     case 2: {
       cout << "Data Penyewa\n";
\end{lstlisting}
\begin{lstlisting}
       for (int i = 1; i < user_count; ++i) {
         cout << i << ". " << "Username: "
            << users[i * 2] << " (Rp."
            << user_bills[i] << ")\n";
       }
       app_state = 1; pause();
     } break; } }
\end{lstlisting}
    \ilz{}{\Script} diatas merupakan sebuah fungsi yang tidak mengembalikan niali apapun atau fungsi \inputscript{void}. Fungsi tersebut akan menangani semua fungsi yang berkaitan dengan interaksi pengguna akun admin dengan program.

\begin{lstlisting}
 int main() {
   while (true) {
     cout << "\33[H\33[2J";
     if (app_state == 0x6) return 0;
     else if (app_state == 0x9) {
       user_logged_id = -1;
       app_state = 0;
     }
     else if (user_logged_id == 0 &&
          app_state == 0x7) {
       user_logged_id = -1;
       app_state = 0;
     }
     switch (app_state & 0x1) {
       case 0: {
         unlogged_page();
       } break;
       case 1: {
         if (user_logged_id == 0) admin_page();
         else logged_page();
       } } } }
\end{lstlisting}
    \ilz{}{\Script} diatas merupakan \textit{entry point} dari program Rental PS Rijal. Fungsi tersebut akan menjalankan fungsi-fungsi lainnya tergantung dengan nilai \inputscript{app\_state}. Program tersebut akan dijalankan selamanya, selama nilai \inputscript{app\_state} bukan 9.

    \subsection{Permasalahan Matematika Diskrit dan Aljabar Linier}

\begin{lstlisting}
 int factorial(int n, int acc = 1)
 {
   if (n == 1) return acc;
   return factorial(n - 1, acc * n);
 }
\end{lstlisting}
    \ilz{}{\Script} diatas mengandung sebuah fungsi yang mengembalikan nilai bilangan bulat dengan mengambil masukan berupa \inputscript{int n}. Fungsi tersebut akan mengembalikan sebuah nilai faktorial dari parameter yang dimasukan dengan menggunakan metode rekursif.

\begin{lstlisting}
   cout << "1. Permutasi\n";
   cout << "2. Kombinasi\n";
   cout << "3. SPL\n";
   int choice;
\end{lstlisting}
\begin{lstlisting}
   while (true) {
     cout << "Pilih: ";
     cin >> choice;
     if (cin.fail() || choice < 1
       || choice > 3) {
       cin.clear();
       cin.ignore();
       cout << "Masukan tidak valid." << endl;
       continue;
     } break;
   }
\end{lstlisting}
    \ilz{}{\Script} diatas akan mencetak tiga pilihan yang diberikan kepada pengguna. Ketiga pilihan tersebut ialah permutasi, kombinasi dan sistem persamaan linier. Masukan akan diminta kepada pengguna kemudian akan dicek validitasnya dengan cara menanyakan apakah masukan lebih kecil dari 1 atau lebih besar dari 3.

\begin{lstlisting}
     case 1: {
       int n, r;
       cout << "Masukan n: ";
       cin >> n;
       cout << "Masukan r: ";
       cin >> r;
       cout << "Hasil Permutasi: " 
          << factorial(n)/factorial(n - r)
          << endl;
     } break;
\end{lstlisting}
    \ilz{}{\Script} diatas digunakan untuk menghasilkan nilai permutasi dari dua masukan. Nilai permutasi dicari dengan penggunakan fungsi \inputscript{factorial} dan meletakannya sesuai rumus permutasi.

\begin{lstlisting}
     case 2: {
       int n, r;
       cout << "Masukan n: ";
       cin >> n;
       cout << "Masukan r: ";
       cin >> r;
       cout << "Hasil Kombinasi: " 
          << factorial(n)/(factorial(n - r) * factorial(r))
          << endl;
     } break;
\end{lstlisting}
    \ilz{}{\Script} diatas digunakan untuk menghasilkan nilai kombinasi dari dua masukan. Nilai kombinasi dicari dengan penggunakan fungsi \inputscript{factorial} dan meletakannya sesuai rumus kombinasi.

\begin{lstlisting}
    case 3: {
      cout << "Masukan Nilai Matriks 3x3\n";
      float matrix[12];
      for (int i = 0; i < 3; ++i) {
        for (int j = 0, mid = 0; j < 3; ++j) {
          cin >> matrix[i * 4 + j];
\end{lstlisting}
\begin{lstlisting}
        }
      }
      cout << "Masukan Nilai B\n";
      for (int i = 0; i < 3; ++i) {
        cin >> matrix[(i + 1) * 4 - 1];
      } cout << "Augmented Matriks\n";
      for (int i = 0; i < 3; ++i) {
        for (int j = 0; j < 4; ++j) {
          cout << matrix[i * 4 + j] << " ";
        } cout << endl;
      }
      float x;
      if (matrix[4] != 0 && matrix[8] != 0
        && matrix[9] != 0) {
        x = matrix[4]/matrix[0];
        for (int i = 0; i < 4; ++i)
          matrix[4 + i] += matrix[i] * -x;
        x = matrix[8]/matrix[0];
        for (int i = 0; i < 4; ++i)
          matrix[8 + i] += matrix[i] * -x;
        x = matrix[9]/matrix[5];
        for (int i = 0; i < 4; ++i)
          matrix[8 + i] += matrix[4 + i] * -x;
      }
      cout << "Matriks Segitiga Atas | SPL\n";
      for (int i = 0; i < 3; ++i) {
        for (int j = 0; j < 3; ++j) {
          cout << matrix[i * 4 + j] << " ";
        } cout << " | " << matrix[(i + 1) * 4 - 1] << endl;
      }
      float z = matrix[11]/matrix[10];
      float y = (matrix[7] - (matrix[6] * z))/matrix[5];
      x = (matrix[3] - (matrix[2] * z) - matrix[1] * y)/matrix[0];
      cout << "x: " << x << "\n";
      cout << "y: " << y << "\n";
      cout << "z: " << z << "\n";
    } break; } }
\end{lstlisting}
    \ilz{}{\Script} diatas digunakan untuk menghasilkan sistem persamaan linier yang dihasilkan dari masukan berupa matriks 3 kali 3 dan vektor B yang digabungkan menjadi satu buah matriks 3 kali 4. Matriks tersebut akan dibuat menjadi matriks segitiga atas yang kemudian akan diproses menjadi sistem persamaan linier.
\end{document}
